% Round 12. Owner: r12_alteration_descent.
% This section gives algebraic descent of numerical witnesses, then applies
% the actual model-Chern and exact-support results of Rounds 10 and 11.
\section{Descent of numerical witnesses and actual Green forms}
\label{r12a:section}

\paragraph{Scope and provenance.}
The supported algebraic projection argument below was already isolated
in the Round 9 foundations audit.  We give its proof, including the
support on which a finite-field degree is taken, and connect it to the
actual model-Chern detector and exact-support Green-pair construction.
The resulting implication is stronger than the finite good-cover
criterion of Round 11: a numerical witness on a proper alteration is
enough when a regular model of the original variety is supplied.
It does not construct a proper pushforward of analytic forms.

Fix the original finite extension $K/\mathbf Q_p$, put $R=K^\circ$,
and write $k$ for its residue field.  A regular model below is integral,
regular, projective and flat over $R$, with a smooth integral generic
fiber.  Models over the integers $R'$ of a finite extension $K'/K$
may be regarded as $R$-models.  Their scheme-theoretic special fibers,
including all multiplicities, are retained.  Degrees on a vertical
proper scheme are degrees over $k$ unless a larger residue field is
explicitly specified.  All Chow groups and $K$-groups are rational.
Every actual Green-form and model-curvature-obstruction conclusion below
is in the range $1\le r\le\dim X$.  The smooth integral generic fibers
need not be geometrically integral over $K$: the exact-support and global
primitive constructions have this native integral scope by
Lemma~\ref{r12i:integral-PC}.

\subsection{Pure operational representatives and supported projection}

For a regular model $\mathcal X$ and a homogeneous class
$\Theta\in\operatorname{CH}^r(\mathcal X)_{\mathbf Q}$, denote by
$\xi_\Theta\in K_0(\operatorname{Vect}(\mathcal X))_{\mathbf Q}$
the unique class with
\[
                 \operatorname{ch}(\xi_\Theta)\cap[\mathcal X]
                              =\Theta.
\tag{ND1}\label{r12a:pure-preimage}
\]
Its existence and uniqueness follow from Stacks, Section 42.58,
Lemma 42.58.1 (Tags 0FEX and 0FEY).  This statement applies to the
complete DVR base; it does not require the model to be smooth over $R$.
The operation
\[
                  c_\Theta:=\operatorname{ch}_r(\xi_\Theta)
                          \in A^r(\mathcal X)_{\mathbf Q}
\]
is a genuine operational class.  In fact all other components of
$\operatorname{ch}(\xi_\Theta)$ vanish operationally.  To see this,
apply $2^{-r}\psi^2$ to \eqref{r12a:pure-preimage}.  It gives the same
Chow class, so uniqueness gives
$\psi^2\xi_\Theta=2^r\xi_\Theta$.  The Adams identity then gives
$(2^j-2^r)\operatorname{ch}_j(\xi_\Theta)=0$ for $j\ne r$.
This is the purity argument immediately following Lemma 42.58.1.
Thus purity is available on every vertical support, not just after
capping with $[\mathcal X]$.

Call $\Theta$ \emph{vertically numerically zero} if
\[
          \deg_k\bigl(c_\Theta\cap[V]\bigr)=0
       \quad\text{for every integral }r\text{-dimensional }
                    V\subset\mathcal X_s.
\tag{ND2}\label{r12a:vertical-zero}
\]
The zero-cycle class inside the degree is a class on the proper
$k$-scheme $V$.  We never take an unweighted degree on
$\operatorname{CH}_0(\mathcal X)$.

\begin{lemma}[Supported numerical pushdown]
\label{r12a:supported-pushdown}
Let $f:\mathcal Y\to\mathcal X$ be a projective morphism between
regular models, of dimensions $M$ and $N$.  Put $\delta=M-N$.
For $\Theta\in\operatorname{CH}^p(\mathcal Y)_{\mathbf Q}$, put
$q=p-\delta$ and assume $0\le q\le N-1$.
For every integral $q$-dimensional $V\subset\mathcal X_s$, there is
an identity
\[
 c_{f_*\Theta}\cap[V]
   =f_{V*}\bigl(c_\Theta\cap f^![V]\bigr)
                         \quad\text{in }\operatorname{CH}_0(V)_{\mathbf Q},
\tag{ND3}\label{r12a:supported-projection}
\]
where $f_V:f^{-1}V\to V$ and the refined Gysin class $f^![V]$ belongs
to $\operatorname{CH}_p(f^{-1}V)_{\mathbf Q}$.
Consequently $f_*\Theta$ is vertically numerically zero whenever
$\Theta$ is.
\end{lemma}

\begin{proof}
A morphism between these regular schemes is lci.  Because $f$ is
projective, its factorization through a projective bundle gives the
factorable Gysin operation $f^!$, of virtual relative dimension
$\delta$, with $f^![\mathcal X]=[\mathcal Y]$.
These are the hypotheses used in Stacks, Lemma 42.59.13 (Tag 0FFA).

The coherent sheaf $\mathcal O_V$ is perfect on the regular ambient
$\mathcal X$.  Its localized leading character supplies a supported
bivariant operation
\[
 \lambda_V=
  \operatorname{ch}^{V}_{N-q}(\mathcal O_V)
          \in A^{N-q}(V\longrightarrow\mathcal X)_{\mathbf Q},
              \qquad \lambda_V\cap[\mathcal X]=[V].
\]
The last equality holds on $V$, by Stacks, Proposition 42.57.1
(Tag 0FEW).  Neither $V$ nor its inverse image is required to be regular.

Ordinary Chern operations are central.  Proper pushforward commutes
with a bivariant operation, and a factorable lci Gysin operation
commutes with an arbitrary supported bivariant operation by Stacks,
Lemma 42.59.7 (Tag 0FF6).  Applying these facts gives
\[
\begin{aligned}
 c_{f_*\Theta}\cap[V]
 &=\lambda_V\cap\bigl(c_{f_*\Theta}\cap[\mathcal X]\bigr)\\
 &=\lambda_V\cap f_*\Theta\\
 &=f_{V*}\bigl((f^*\lambda_V)\cap\Theta\bigr)\\
 &=f_{V*}\bigl(c_\Theta\cap
               ((f^*\lambda_V)\cap f^![\mathcal X])\bigr)\\
 &=f_{V*}\bigl(c_\Theta\cap
                          f^!(\lambda_V\cap[\mathcal X])\bigr)\\
 &=f_{V*}\bigl(c_\Theta\cap f^![V]\bigr).
\end{aligned}
\]
Every term is supported on $V$ or $f^{-1}V$.  In particular this
argument proves \eqref{r12a:supported-projection}, which is stronger
than the projection formula after forgetting the support.

The refined class $f^![V]$ has dimension $q+\delta=p$ and is supported
on the special fiber.  Choose a cycle representative, a finite rational
sum of integral vertical $p$-cycles.  The degrees after applying
$c_\Theta$ vanish termwise by \eqref{r12a:vertical-zero}.  Taking
degree in \eqref{r12a:supported-projection} proves the assertion.
\end{proof}

\begin{lemma}[Numerical pullback]
\label{r12a:numerical-pullback}
For a projective morphism $f:\mathcal Y\to\mathcal X$ between regular
models, the class $f^!\Theta$ has operational representative
$f^*c_\Theta$.  If $\Theta$ is vertically numerically zero, then so
is $f^!\Theta$.
\end{lemma}

\begin{proof}
The class $f^*\xi_\Theta$ has pure operational Chern character
$f^*c_\Theta$.  Chern operations commute with $f^!$, hence
\[
 (f^*c_\Theta)\cap[\mathcal Y]
       = f^!(c_\Theta\cap[\mathcal X])=f^!\Theta.
\]
Uniqueness in \eqref{r12a:pure-preimage} proves the first assertion.
For an integral vertical $p$-cycle $W$ on $\mathcal Y$, with
$p=\operatorname{codim}\Theta$, proper operational projection gives
\[
 \deg_k\bigl((f^*c_\Theta)\cap[W]\bigr)
       =\deg_k\bigl(c_\Theta\cap f_*[W]\bigr)=0.
\]
Here $f_*[W]$ is a vertical $p$-cycle or zero.  This proves the second
assertion, without flatness of the morphism on special fibers.
\end{proof}

\paragraph{Constant fields.}
If the source is originally a model over $R'$ with residue field $k'$,
vertical numerical vanishing computed over $k'$ implies the same
vanishing over $k$: the two degrees differ by $[k':k]$.  Refined Gysin
classes and proper pushforwards retain their scheme-theoretic
multiplicities.  In particular no ramification multiplicity is replaced
by one when the model is viewed over $R$.

\subsection{The alteration witness descends to a regular original model}

\begin{theorem}[A regular original model accepts proper numerical witnesses]
\label{r12a:regular-target}
Let $X/K$ be smooth integral projective, and suppose a regular
projective $R$-model $\mathcal X$ of $X$ is supplied.  Let
$f:\mathcal Y\to\mathcal X$ be a projective dominant generically finite
morphism from a regular model, with generic map $h:Y\to X$ of native
degree $d=[K(Y):K(X)]$.  Let $a\in\operatorname{CH}^r(X)_{\mathbf Q}$,
where $1\le r\le\dim X$.
If there is a vertically numerically zero
$\Theta\in\operatorname{CH}^r(\mathcal Y)_{\mathbf Q}$ with
\[
                              \Theta|_Y=h^*a,
\tag{ND4}\label{r12a:witness-generic}
\]
then $\mathfrak o_X^r(a)=0$.  Every prescribed cycle representing
$a$ has a strong zero-curvature actual Green form over the original
$K$, on its exact complement, against all actual tests.
\end{theorem}

\begin{proof}
Set
\[
                           \Lambda=\frac1d f_*\Theta
                      \in\operatorname{CH}^r(\mathcal X)_{\mathbf Q}.
\tag{ND5}\label{r12a:descended-class}
\]
The generic projection formula on the smooth varieties gives
\[
          \Lambda|_X=\frac1d h_*h^*a=a.
\]
By Lemma~\ref{r12a:supported-pushdown}, $\Lambda$ is vertically
numerically zero.  Choose its pure virtual vector-bundle representative
$\xi_\Lambda$ from \eqref{r12a:pure-preimage}.
Then
\[
 \operatorname{ch}_r(\xi_\Lambda|_X)=a,
 \qquad
 \deg_k\bigl(\operatorname{ch}_r(\xi_\Lambda|_{\mathcal X_s})
                                          \cap[V]\bigr)=0
\]
for every integral vertical $r$-cycle $V$.

The Round 10 numerical model-Chern detector gives the strict actual
identity $\operatorname{ch}_r(\bar\xi_\Lambda)=0$ on $X^{\rm an}$.
Round 11's well-defined model-curvature class gives
$\mathfrak o_X^r(a)=0$, and its exact-support Green-pair theorem then
gives the stated zero-curvature form for each prescribed representative.
The virtual class need not be supported on that representative: the
global generic-Chow-zero correction in Round 11 is exactly what removes
that unnecessary condition.

Only algebraic classes were pushed through $f$.  The analytic form is
constructed on the original $X$ after \eqref{r12a:descended-class};
there is no proper pushforward of a Green form or a primitive.
\end{proof}

\begin{corollary}[A supplied witness need not map to the chosen model]
\label{r12a:generic-witness}
In Theorem~\ref{r12a:regular-target}, it suffices to have a proper
generically finite morphism $h:Y\to X$ of smooth projective generic
fibers, a regular projective model $\mathcal Y$ of $Y$ over a finite
extension of $R$, and a vertically numerically zero lift of $h^*a$
on $\mathcal Y$.  A model morphism to $\mathcal X$ is not required
as part of the supplied data.
\end{corollary}

\begin{proof}
Take the closure of the generic graph of $h$ in
$\mathcal Y\times_R\mathcal X$.  It is integral, projective and flat
over $R$, and maps to both models.  A regular projective alteration
$\mathcal W$ of this closure exists by the projective case of
de Jong's Theorem 6.5; its constant field may be a further finite
extension.  Regard all schemes as $R$-schemes.
Write $b:\mathcal W\to\mathcal Y$ and
$g:\mathcal W\to\mathcal X$ for the resulting maps.
They are projective between regular schemes, and
$b^!\Theta$ is vertically numerically zero by
Lemma~\ref{r12a:numerical-pullback}.  On generic fibers its restriction
is $g_K^*a$.  Apply Theorem~\ref{r12a:regular-target} to $g$.
The degree used there is the total function-field degree of $g$,
including its constant-field factor.
\end{proof}

Thus the numerical-witness property is independent of which supplied
regular projective original model is chosen: use a common regular
alteration and the two preceding lemmas.  This is a statement about
numerical witnesses.  It does not establish birational invariance or
pullback injectivity for the whole actual Bott--Chern quotient.

\subsection{Finite numerical covers without an original regular model}

\begin{theorem}[A finite numerical-model cover]
\label{r12a:finite-numerical-cover}
Let $h:Y\to X$ be finite dominant between smooth integral projective
$K$-varieties, and let $a\in\operatorname{CH}^r(X)_{\mathbf Q}$,
where $1\le r\le\dim X$.
Suppose an integral projective flat model $\mathcal Y$ of $Y$ carries
a virtual vector bundle $\xi$ with
\[
 \operatorname{ch}_r(\xi|_Y)=h^*a,
 \qquad
 \deg_k\bigl(\operatorname{ch}_r(\xi|_{\mathcal Y_s})\cap[V]\bigr)=0
\tag{ND6}\label{r12a:finite-numerical-data}
\]
for every integral vertical $r$-cycle $V$.
Then $\mathfrak o_X^r(a)=0$ and every prescribed cycle in class $a$
satisfies strong G on its exact complement, over $K$, on all actual
tests.  No regular model of $X$ is assumed.

The same statement holds when $Y$ is smooth integral over a finite
extension $K'/K$, its model is over $R'$, and $h$ is a finite dominant
$K$-morphism.  In that case the vanishing condition may be checked
using degrees over the residue field of $R'$.
\end{theorem}

\begin{proof}
The numerical detector gives
$\Omega_Y:=\operatorname{ch}_r(\bar\xi)=0$ as a $K$-actual form.
For the extension-field version, view the $R'$-model over $R$;
every integral cycle on the thickened special fiber lies on its
reduction, and the degree changes by the residue-field degree.
Thus this conclusion is again obtained directly in the original
$K$-actual theory, without comparing test spaces over different fields.
Lemma~\ref{r12i:integral-PC} supplies the exact-support and global
primitive constructions in this smooth integral, possibly
nongeometrically integral, $K$-scope.

Choose a model character $\Omega_X$ representing $a$.  On a common
model of $Y$, the character $h^*\Omega_X-\Omega_Y=h^*\Omega_X$ has
generic degree-$r$ Chow class zero.  The constructive proof of Round
11's generic-Chow-zero primitive theorem writes a primitive as
\[
\begin{gathered}
 \beta_Y=\sum_j q_j e_j^{-1}
                  \operatorname{Tr}^{\mathscr M}_{b_j}(G_j),\\
 DG_j=b_j^*\Omega_j,\qquad
 \sum_j q_j\Omega_j=h^*\Omega_X.
\end{gathered}
\tag{ND7}\label{r12a:primitive-generators}
\]
Here each $b_j:T_j\to Y$ is finite flat on all of $Y$, of degree
$e_j$, and each $G_j$ is an original global sum of PL scalars times
model curvature polynomials.  These properties follow because the
complexes used for a generic $K_0$ relation are acyclic on the full
generic fiber.  Adams projection handles a relation in only the
specified Chow degree.

Since $h$ is finite between smooth varieties of the same dimension,
it is flat.  Put $d=\deg_K h$ and define directly
\[
 \beta_X=\frac1d\sum_jq_j e_j^{-1}
                 \operatorname{Tr}^{\mathscr M}_{h\circ b_j}(G_j).
\tag{ND8}\label{r12a:finite-primitive}
\]
All traces are permitted by the Round 11 minimal-trace theorem.
We never trace the possibly larger-class form $\beta_Y$ itself.
For every compact complementary actual test $\eta$, trace pairing,
actual Stokes, and the native degree formula give
\[
\begin{aligned}
 \int_X\beta_XD\eta
 &=\frac1d\sum_jq_j e_j^{-1}
                  \int_{T_j}G_j(hb_j)^*D\eta\\
 &=\frac1d\sum_jq_j\int_Y\Omega_jh^*\eta\\
 &=\frac1d\int_Yh^*(\Omega_X\eta)
  =\int_X\Omega_X\eta.
\end{aligned}
\]
Current injectivity gives $D\beta_X=\Omega_X$ as actual forms.
Hence $\mathfrak o_X^r(a)=0$, and the exact-support theorem proves G.
\end{proof}

\begin{corollary}[A finite cover with a regular model]
\label{r12a:finite-regular-cover}
Let $h:Y\to X$ be as in Theorem~\ref{r12a:finite-numerical-cover}.
Suppose $Y$ has a supplied regular projective model.  If $h^*a$ has
a vertically numerically zero lift on a regular projective proper
alteration of $Y$, then $\mathfrak o_X^r(a)=0$ and strong G holds for
every prescribed representative of $a$.

In particular, a supplied regular model after a finite constant-field
extension is enough for this implication, provided the stated numerical
witness has been established on an alteration.
\end{corollary}

\begin{proof}
The proof of Theorem~\ref{r12a:regular-target} and
Corollary~\ref{r12a:generic-witness} gives an actual virtual model
class $\xi_\Lambda$ on the supplied regular model of $Y$, with the
two properties in \eqref{r12a:finite-numerical-data}.
Apply Theorem~\ref{r12a:finite-numerical-cover}.
\end{proof}

\subsection{Why an arbitrary alteration is still different}

\begin{proposition}[No finite dominating cover of a nontrivial smooth blowup]
\label{r12a:no-finite-blowup-cover}
Let $X$ be a smooth integral variety over a field, and let
$C\subset X$ be a nonempty smooth closed subscheme of codimension
at least two.  Put $p:Y=\operatorname{Bl}_C X\to X$.
There is no finite dominant morphism $q:T\to X$ from an integral
variety that factors as $q=p\circ u$ for a morphism $u:T\to Y$.
\end{proposition}

\begin{proof}
Let $E=p^{-1}C$ be the exceptional effective Cartier divisor.
The blowup property is Stacks, Lemmas 31.33.4--31.33.5
(in particular Tag 0806).  If such a factorization existed,
$u$ would be dominant: over $X\setminus C$, the map $p$ is an
isomorphism and the image of $q$ is dense.
Since $T$ is integral and $u$ is dominant, a local equation of $E$
pulls back to a nonzero element of a domain.  Thus $u^{-1}E$ is an
effective Cartier divisor.  It is nonempty, since the finite dominant
$q$ is surjective.  It has dimension $\dim T-1=\dim X-1$.

On the other hand, scheme-theoretically $u^{-1}E=q^{-1}C$.
The latter is finite over $C$, so has dimension at most
$\dim C\le\dim X-2$.  This contradiction proves the claim.
\end{proof}

The proposition rules out a direct replacement of a general alteration
by a finite multisection over the original target.  The finite
multisection in the supported-complex construction works because its
chosen smooth flag fibration has constant positive fiber dimension;
the blowup has no such property.  Stein factorization does not repair
the problem either: the finite factor of a proper birational morphism
to a normal target is the target itself, leaving the whole modification
in the connected-fiber factor.

Theorems~\ref{r12a:regular-target} and
\ref{r12a:finite-numerical-cover} avoid this obstruction by two distinct
verified mechanisms: algebraic supported pushdown to a regular model,
or analytic minimal trace along an actually finite generic morphism.
Neither mechanism proves descent from an arbitrary proper alteration
when no regular model of the original variety, or of a finite smooth
cover, is supplied.  Nor does a merely existing primitive upstairs
automatically have the special generator presentation
\eqref{r12a:primitive-generators}.

\paragraph{Primary sources checked for this section.}
\begin{itemize}
\item Stacks, rational intersection on regular schemes and the pure
operational $K_0$ representative:
\url{https://stacks.math.columbia.edu/tag/0FEX}.
\item Stacks, supported leading Chern character:
\url{https://stacks.math.columbia.edu/tag/0FEW}.
\item Stacks, arbitrary supported bivariant commutation with lci Gysin,
and the regular-scheme projection formula:
\url{https://stacks.math.columbia.edu/tag/0FF6},
\url{https://stacks.math.columbia.edu/tag/0FFA}.
\item A. J. de Jong, \emph{Smoothness, semi-stability and alterations},
Theorem 6.5, printed p.83:
\url{https://www.numdam.org/item/PMIHES_1996__83__51_0.pdf}.
For a projective input, properness makes the open projective
compactification in the theorem equal to its source.
\item Stacks, the exceptional Cartier divisor and blowup universal
property: \url{https://stacks.math.columbia.edu/tag/0806}.
\end{itemize}
